% ============================================================
%  Chapitre 2 — Compréhension métier
% ============================================================
\chapter{Compréhension métier}
\label{ch:metier}

\section{Introduction}

Ce chapitre formalise le besoin auquel le système doit répondre. Nous
caractérisons d'abord la problématique de l'accès à l'information documentaire
en entreprise, puis nous analysons les solutions existantes et leurs limites.
Nous exposons enfin la solution proposée et les critères objectifs retenus pour
juger de son succès.

\section{Problématique}

Une entreprise de taille moyenne accumule, sur quelques années d'activité,
plusieurs dizaines de milliers de documents hétérogènes. Cette accumulation
crée un paradoxe : l'information est présente et conservée, mais son coût
d'accès la rend pratiquement inutilisable. Trois obstacles distincts se
combinent.

\subsection{Le fossé lexical}

Les moteurs de recherche documentaires classiques reposent sur la
correspondance de mots-clés. Une requête ne trouve un document que si les termes
employés par l'utilisateur figurent littéralement dans ce document. Or les
rédacteurs et les lecteurs n'emploient pas le même vocabulaire : une question
portant sur les « congés » ne remontera pas une procédure intitulée « absences
autorisées », et une recherche sur « fournisseur » ignorera un document parlant
de « prestataire ». Ce phénomène, connu sous le nom de \emph{fossé lexical}
(\eng{vocabulary mismatch}), constitue la limite structurelle de la recherche
par mots-clés.

\subsection{La charge de synthèse}

Même lorsqu'un moteur classique retrouve les bons documents, il retourne une
\emph{liste de documents}, non une \emph{réponse}. L'utilisateur doit encore
ouvrir chaque résultat, localiser le passage utile et synthétiser mentalement.
Pour une question dont la réponse est répartie sur plusieurs documents, ce
travail est long et sujet à l'oubli.

\subsection{L'hallucination des modèles génératifs}

Le recours direct à un modèle de langue généraliste, tel que ChatGPT, semble
apporter une réponse immédiate à ces deux obstacles : il comprend le langage
naturel et produit une réponse rédigée. Deux difficultés majeures l'interdisent
cependant en contexte professionnel.

D'abord, le modèle \textbf{ne connaît pas les documents de l'entreprise}. Ils
n'appartiennent pas à ses données d'entraînement, et ils ne le peuvent pas :
transmettre des contrats, des factures ou des données salariales à un service
tiers pose un problème de confidentialité et de conformité réglementaire
(RGPD) rédhibitoire pour la plupart des organisations.

Ensuite, et plus gravement, le modèle \textbf{hallucine}. Lorsqu'il ne dispose
pas de l'information demandée, il ne répond pas « je ne sais pas » : il génère
la suite de texte statistiquement la plus vraisemblable, c'est-à-dire une
réponse plausible dans sa forme mais fausse dans son contenu. Un assistant
interne qui inventerait un montant de facture, une date d'échéance contractuelle
ou une étape de procédure de sécurité serait pire qu'inutile : il serait
dangereux, d'autant que sa formulation assurée n'offre à l'utilisateur aucun
signal d'alerte.

\begin{encadre}
\textbf{Formulation de la problématique.} Comment permettre à un collaborateur
d'interroger en langage naturel l'ensemble des documents internes de son
entreprise, en obtenant une réponse directe, \emph{vérifiable} et
\emph{traçable}, sans que ces documents ne quittent l'infrastructure de
l'organisation ?
\end{encadre}

\section{Étude de l'existant}

Plusieurs familles de solutions adressent tout ou partie de cette
problématique. Le tableau~\ref{tab:existant} en présente une analyse comparative.

\begin{table}[htbp]
\centering
\caption{Analyse comparative des solutions existantes}
\label{tab:existant}
\small
\begin{tabularx}{\textwidth}{@{}p{3.1cm}XX@{}}
\toprule
\textbf{Solution} & \textbf{Avantages} & \textbf{Inconvénients} \\
\midrule

Recherche plein texte (Elasticsearch, Solr) &
Mature, très rapide, passe à l'échelle sur des millions de documents ; aucun
coût d'inférence &
Sensible au fossé lexical ; retourne des documents et non des réponses ;
la charge de synthèse reste entièrement sur l'utilisateur \\

\addlinespace
LLM généraliste seul (ChatGPT, GPT-4o) &
Excellente compréhension du langage naturel ; réponse directe et bien rédigée ;
aucune infrastructure à maintenir &
N'a aucune connaissance des documents internes ; hallucine sans le signaler ;
envoi de données confidentielles à un tiers ; coût d'API récurrent \\

\addlinespace
Suites propriétaires (Microsoft 365 Copilot, Glean) &
Intégration native à l'écosystème bureautique ; RAG déjà implémenté ;
maintenance déléguée &
Coût par utilisateur élevé ; suppose que les documents soient déjà hébergés chez
l'éditeur ; fonctionnement opaque, métriques de fiabilité non exposées ;
dépendance forte au fournisseur \\

\addlinespace
Assistants de base de connaissances (Notion AI, Confluence AI) &
Simples d'usage ; bien intégrés à l'outil hôte &
Limités aux documents de la plateforme ; ne couvrent pas les PDF et archives
externes ; traçabilité des sources partielle \\

\addlinespace
\textbf{RAG sur mesure} (LangChain, ChromaDB, Ollama) &
\textbf{Fonctionne entièrement en local} ; aucune donnée ne sort de
l'infrastructure ; aucun coût d'API ; sources citées et vérifiables ; métriques
de fiabilité mesurables ; contrôle total sur chaque composant &
Nécessite un développement et une maintenance internes ; latence supérieure sur
matériel sans GPU ; qualité dépendante du réglage du pipeline \\

\bottomrule
\end{tabularx}
\end{table}

Cette analyse fait ressortir un constat : aucune solution du marché ne réunit
simultanément la confidentialité totale, l'absence de coût récurrent et la
traçabilité vérifiable des réponses. C'est cet espace que le projet vise à
occuper.

\section{Solution proposée}

Nous proposons de construire un assistant documentaire interne fondé sur
l'architecture \textbf{RAG} (\eng{Retrieval-Augmented Generation}), dont le
principe est de découpler la \emph{connaissance} de la \emph{formulation}.

La connaissance reste dans les documents de l'entreprise, indexés dans une base
vectorielle locale. Lorsqu'une question est posée, le système commence par
rechercher les passages documentaires les plus proches sémantiquement de cette
question. Ces passages sont ensuite injectés dans l'invite envoyée au modèle de
langue, accompagnés d'une consigne stricte : répondre uniquement à partir du
contexte fourni, et déclarer explicitement son ignorance si la réponse ne s'y
trouve pas. Le modèle n'est plus une source de connaissance, mais un moteur de
reformulation contraint.

Ce basculement change la nature du problème de l'hallucination. Le modèle ne
peut plus inventer à partir de son savoir général, puisque l'invite lui interdit
de sortir du contexte fourni ; et lorsqu'il le fait malgré tout, la présence des
extraits sources permet à l'utilisateur de le détecter immédiatement.

Les choix technologiques découlent directement de la contrainte de
confidentialité :

\begin{itemize}
  \item \textbf{ChromaDB} comme base vectorielle, persistée sur disque local,
        plutôt qu'un service géré comme Pinecone ;
  \item \textbf{sentence-transformers} pour la vectorisation, exécuté en local
        sur CPU, plutôt que l'API \texttt{text-embedding-ada-002} d'OpenAI ;
  \item \textbf{Ollama} avec le modèle \textbf{Mistral 7B} pour la génération,
        plutôt que GPT-4o ;
  \item \textbf{LangChain} pour l'orchestration de la chaîne ;
  \item \textbf{RAGAS} pour la mesure de la qualité et du taux d'hallucination.
\end{itemize}

Aucun de ces composants n'émet de requête vers un service externe en phase
d'exploitation. L'ensemble de la chaîne, de l'ingestion à la génération,
s'exécute sur la machine hôte.

\section{Critères de succès}

Conformément à la première phase de CRISP-DM, des critères objectifs ont été
définis en amont pour juger de la réussite du projet. Ils sont présentés au
tableau~\ref{tab:criteres} et repris au chapitre~\ref{ch:evaluation}.

\begin{table}[htbp]
\centering
\caption{Critères de succès définis en phase de compréhension métier}
\label{tab:criteres}
\small
\begin{tabularx}{\textwidth}{@{}p{4.2cm}Xp{2.9cm}@{}}
\toprule
\textbf{Critère} & \textbf{Description} & \textbf{Mesure} \\
\midrule
Ancrage documentaire &
Toute affirmation de la réponse doit être soutenue par le contexte récupéré &
\eng{Faithfulness} (RAGAS) \\
Pertinence de la réponse &
La réponse doit adresser la question effectivement posée &
\eng{Answer relevancy} \\
Qualité de la recherche &
Les extraits récupérés doivent être pertinents et complets &
\eng{Context precision} et \eng{recall} \\
Traçabilité &
Chaque réponse doit exposer ses documents sources et leur score &
Vérification fonctionnelle \\
Confidentialité &
Aucune donnée ne doit quitter l'infrastructure &
Audit des dépendances réseau \\
Aveu d'ignorance &
Le système doit refuser de répondre hors contexte &
Test sur questions hors corpus \\
\bottomrule
\end{tabularx}
\end{table}

\section{Conclusion}

Ce chapitre a formalisé la problématique du projet en identifiant ses trois
composantes : le fossé lexical qui limite la recherche par mots-clés, la charge
de synthèse laissée à l'utilisateur, et le risque d'hallucination des modèles
génératifs employés seuls. L'étude de l'existant a montré qu'aucune solution du
marché ne concilie confidentialité, gratuité et traçabilité. La solution
proposée --- une architecture RAG entièrement locale --- répond à ces trois
exigences. Le chapitre suivant s'intéresse au corpus documentaire sur lequel ce
système a été construit.
